\scanpage{40}
\noindent (23)
\diagram{\begin{tikzcd}
\mathbb Z\arrow[r]\arrow[d]&1\\
1
\end{tikzcd}},
on a maintenant une équivalence de topos
\[
\deleted{\unsure{(24)}}\qquad \unsure{\mathcal F|\mathfrak F}\unsure{\cong|\simeq}\widehat C,
\]
où $C$ est la catégorie \diagram{\begin{tikzcd}[column sep=small]\bullet\arrow[r]&\bullet\end{tikzcd}}, qui a
un objet final --- donc $\unsure{\mathcal F|\mathfrak F}$ \uline{est encore}
\uline{homotopiquement trivial,} \uline{alors que cette}
\uline{fois-ci $X$ ne l'est pas} ! Pourtant les
pièces de construction $\{a\}$, $X^*\simeq\mathbb C$, et
$\unsure{V}^*_{z,Z}\simeq\mathbb D^*_{\infty}$ sont bel et bien des $K(\pi,1)$ ---
(les deux premières $=$ contractibles !). Cela
montre que le topos $\unsure{\mathcal F|\mathfrak F}$, \deleted{\unsure{où les}}\inserted{qui en}
principe devrait ``tenir compte'' des
$\pi_{\unsure{1|i}}$ des ``pièces de construction'' $Y,X^*,V^*_{Y,X}$,
ne donne pas \inserted{néanmoins} le type d'homotopie
de $X$, $=$ si les ``morceaux'' sont des
$K(\pi,1)$ --- en d'autres termes, le
morphisme canonique
\[
(24)\qquad X\longrightarrow\unsure{\mathcal F|\mathfrak F}
\]
n'est pas néc\textsuperscript{t} une équivalence d'homotopie.
\mspage{39}
